\section{赵：从北方边疆到长平的断裂}

赵从\eventref{WQ-0403-THREEJIN}{战国·三家分晋}的遗产中成长。它既有太行以东的农业和人口，也面对北方草原与山地的机动压力。赵的强大并不只有中原贵族政治的一面；它必须不断处理边疆、骑兵、城防、王室继承与中原争霸之间的拉扯。

\begin{event}{公元前307年}{赵武灵王推行胡服骑射}{可靠纪年}{赵}\eventref{WQ-0307-HUFU}{战国·胡服骑射}
\term{这一年发生了什么}：赵武灵王推动服饰与骑射改革，试图训练更适合北方边地战争的机动力量。

\term{事情为什么走到这里}：赵面对林胡、楼烦等北方势力，也要在战国竞争中维持自己的边界。传统礼服与战车体系并非完全无用，但难以解决骑兵机动和边地作战的问题。

\term{史料边界}：《史记·赵世家》保存赵武灵王说服公子成等反对者的大段对话，极具思想戏剧性。它能说明改革触动了身份与礼制象征，不应被当作逐字会议记录。
\end{event}

\begin{event}{公元前295年}{沙丘之乱，赵武灵王死}{可靠纪年}{沙丘}\eventanchor{WQ-0295-SHAQIU}{战国·沙丘之乱}
\term{这一年发生了什么}：赵武灵王在内禅与继承安排失败后被困沙丘，最终死去。

\term{后来改变了什么}：胡服骑射并没有自动产生稳定国家。改革者若不能处理继承、王室成员和军政权力的重新分配，新的军事能力也会被内部危机打断。赵此后仍强，却失去了一位能够将北方战略与中央政治相连的关键人物。
\end{event}

\begin{event}{公元前260年}{长平之战}{可靠纪年}{长平}\eventref{WQ-0260-CHANGPING}{战国·长平之战}
\term{这一年发生了什么}：秦赵长期对峙后，赵军败降。长平不是一场突然爆发的决战，而是围绕上党、补给线和战略选择持续数年的消耗。

\term{事情为什么走到这里}：赵并非没有资源和军事能力，但长期对峙、补给、外交孤立和决策失误共同把边境争夺推成国力消耗战。廉颇被赵括替换也不应只写成一次用人失误，它同时反映赵廷对长期守势和后勤压力的不耐。

\term{后来改变了什么}：赵在这场失败后很难再恢复原有的战略位置。秦东进的门槛降低，东方各国对秦的共同防御也因此更加困难。
\end{event}

\begin{sourcenote}[长平：有尸骨坑，不等于数字已经被验证]
\sourcebook{史记·白起王翦列传}称坑杀赵卒“四十余万”。高平地区发现的战国尸骨坑和兵器遗存，能够支持大规模杀伤事件确实发生；但考古发现的数量不能直接验证“四十余万”这一传世整数。廉颇被赵括替换也不应只写成一次用人失误，还要看到赵廷对长期补给和消耗战的不满。
\end{sourcenote}

\begin{center}
\begin{tikzpicture}[x=.92cm,y=.92cm,font=\small]
  \fill[paper] (0,0) rectangle (10,5.8);
  \fill[dynasty!22] (.6,1.0) -- (3.7,.8) -- (4.3,4.7) -- (1.2,5.0) -- cycle;
  \node[align=center] at (2.35,2.9) {秦\\关中};
  \fill[green!20] (5.3,1.1) -- (9.3,1.0) -- (9.2,5.0) -- (5.7,5.1) -- cycle;
  \node[align=center] at (7.4,3.25) {赵\\邯郸方向};
  \draw[gray!70,line width=1.5pt] (4.7,.3) .. controls (4.2,2.5) and (4.8,4.1) .. (4.5,5.5);
  \node[rotate=90,text=gray!70] at (4.15,3.0) {太行山地};
  \fill[timeline] (5.0,2.9) circle (.13);
  \node[above,align=center] at (5.0,3.05) {长平\\上党一带};
  \draw[-{Stealth[length=2mm]},ultra thick,color=dynasty] (2.8,2.5) -- (4.85,2.9);
  \draw[-{Stealth[length=2mm]},thick,color=green!60!black] (7.3,2.2) -- (5.15,2.85);
  \node[text=dynasty,below] at (3.8,2.3) {秦军东进与补给};
  \node[text=green!60!black,above] at (6.4,2.25) {赵军西向防御};
\end{tikzpicture}

\smallskip
\small\textit{图  长平战场空间示意：突出关中、太行、上党/长平与赵国腹地之间的补给和地形关系；不表示精确行军路线或国界。}
\end{center}


\begin{turningpoint}[制度得失：机动化能解决什么，不能解决什么]
骑射改革提高了赵面对北方的军事适应性，却不能自动解决中原争霸中的粮道、联盟和长期兵源问题。长平提醒读者：军事技术是一项能力，不是一张保证胜利的凭证。
\end{turningpoint}
